\begin{table}[htbp]
\centering
\caption{同一工程认证文档在不同解析方案下的本地对比结果}
\label{tab:local_parser_compare}
\footnotesize
\renewcommand{\arraystretch}{1.2}
\begin{tabular}{p{1.8cm}p{2.5cm}p{1.6cm}p{1.8cm}p{1.7cm}p{1.9cm}p{3.1cm}}
\hline
解析方案 & 测试文档/页面 & 表格结构保留 & 阅读顺序一致性 & 标题层级保留 & 后续切块可用性 & 主要问题 \\
\hline
\texttt{pdftotext -layout} & 培养方案第2页：毕业要求支撑矩阵 & 较差 & 较差 & 一般 & 较差 & 矩阵被线性化展开，行列对应关系难以稳定恢复，不利于按支撑关系组织语义块。 \\
\hline
MinerU pipeline & 培养方案第2页：毕业要求支撑矩阵 & 一般 & 良好 & 良好 & 一般 & 虽可恢复表格外形，但个别“√”被误识别为“专”或波浪号等字符，关键支撑关系存在语义噪声。 \\
\hline
MinerU2.5 & 培养方案第2页：毕业要求支撑矩阵 & 良好 & 良好 & 良好 & 良好 & 表格结构、勾选符号与页内章节边界保持较完整，可直接支撑后续表格块与标题块切分。 \\
\hline
\texttt{pdftotext -layout} & 制度办法第1页与第4页：标题条款及公式表格混排 & 较差 & 较差 & 较差 & 较差 & 对扫描型 PDF 基本无法输出有效文本，第1页与第4页结果为空，无法形成可用语义块。 \\
\hline
MinerU pipeline & 制度办法第4页：表格与公式混排 & 良好 & 一般 & 一般 & 一般 & 表格主体能够恢复，但公式（2）至公式（5）严重失真，标题与条款格式一致性不足。 \\
\hline
MinerU2.5 & 制度办法第1页与第4页：标题条款及公式表格混排 & 良好 & 良好 & 良好 & 良好 & 标题层级、表格结构与公式语义整体保持较好，更适合作为精细化切块与元数据绑定的上游输入。 \\
\hline
\end{tabular}
\end{table}
